% GENERATED FILE. Edit canonical lessons, then run npm run generate:books.
% canonical-source-hash: fnv1a64:2e098fd1d04c230f
% canonical-lessons: FR-C505-arbitre, FR-C505-sifflet, FR-C505-siffler, FR-C505-terrain, FR-C505-marquer, FR-R505-revision-498-505

\chapter{Sur le terrain}
\label{ch:fr-sur-le-terrain}
\hlchaptermodality{505}

\begin{chapteropening}
\textbf{By the end of this chapter:} \emph{I can talk about what happens during a match.}

Complete the last lesson of chapter 505: I can talk about what happens during a match.
\end{chapteropening}

\section[l'arbitre]{l'arbitre — referee, umpire}

\label{lesson:FR-C505-arbitre}



\begin{quote}
Before the new one: say the French for sweat, then the French for to warm up.
\end{quote}

\subsection*{You'll want to know: l'arbitre}
\textbf{l'arbitre} — \textquotedblleft{}referee, umpire\textquotedblright{}.

\begin{quote}
L'arbitre a donné un carton rouge.
\end{quote}

\begin{quote}
\emph{The referee gave a red card.}
\end{quote}

The same form for a man and a woman: \emph{un arbitre}, \emph{une arbitre}.

\begin{cousinweb}
\emph{l'arbitre} is a learned borrowing from Latin \emph{arbiter} \textquotedblleft{}a witness, a judge\textquotedblright{} (EtymDB; Lewis \& Short).
\end{cousinweb}

\begin{grammarlens}[title={What you've built}]
One more word for telling, explaining and describing at B1.
\end{grammarlens}

\subsection*{Your turn}
\begin{itemize}
\raggedright
  \item \emph{Say it:} \emph{l'arbitre}
  \item \emph{Say it:} \emph{l'arbitre}, once more
  \item \emph{Say it:} a sentence of your own with \emph{l'arbitre}
  \item \emph{Recall:} say the French for sweat, then the French for to warm up, then say \emph{l'arbitre} again
\end{itemize}

\begin{tcolorbox}[breakable,colback=teal!4,colframe=teal!35!black,title={Before you move on}]
What does l'arbitre mean? (\textquotedblleft{}Referee, umpire\textquotedblright{}.) Say it once more.
\end{tcolorbox}
\section[le sifflet]{le sifflet — whistle}

\label{lesson:FR-C505-sifflet}



\begin{quote}
Before the new one: say the French for to warm up, then the French for referee.
\end{quote}

\subsection*{You'll want to know: le sifflet}
\textbf{le sifflet} — \textquotedblleft{}whistle\textquotedblright{}.

\begin{quote}
Au coup de sifflet, les joueurs sont sortis.
\end{quote}

\begin{quote}
\emph{At the whistle, the players went off.}
\end{quote}

\emph{Le coup de sifflet final} — the final whistle.

\begin{cousinweb}
\emph{le sifflet} is built on \emph{siffler} \textquotedblleft{}to whistle\textquotedblright{}, which goes back to Late Latin \emph{sīfilāre}, a form of Latin \emph{sībilāre} \textquotedblleft{}to hiss, to whistle\textquotedblright{} (Lewis \& Short).
\end{cousinweb}

\begin{grammarlens}[title={What you've built}]
One more word for telling, explaining and describing at B1.
\end{grammarlens}

\subsection*{Your turn}
\begin{itemize}
\raggedright
  \item \emph{Say it:} \emph{le sifflet}
  \item \emph{Say it:} \emph{le sifflet}, once more
  \item \emph{Say it:} a sentence of your own with \emph{le sifflet}
  \item \emph{Recall:} say the French for to warm up, then the French for referee, then say \emph{le sifflet} again
\end{itemize}

\begin{tcolorbox}[breakable,colback=teal!4,colframe=teal!35!black,title={Before you move on}]
What does le sifflet mean? (\textquotedblleft{}Whistle\textquotedblright{}.) Say it once more.
\end{tcolorbox}
\section[siffler]{siffler — to whistle; to boo}

\label{lesson:FR-C505-siffler}



\begin{quote}
Before the new one: say the French for referee, then the French for whistle.
\end{quote}

\subsection*{You'll want to know: siffler}
\textbf{siffler} — \textquotedblleft{}to whistle; to boo\textquotedblright{}.

\begin{quote}
Le public a sifflé l'équipe après la défaite.
\end{quote}

\begin{quote}
\emph{The crowd booed the team after the defeat.}
\end{quote}

In France, a crowd that whistles is booing.

\begin{cousinweb}
\emph{siffler} is inherited from Late Latin \emph{sīfilāre}, a form of Latin \emph{sībilāre} \textquotedblleft{}to hiss, to whistle\textquotedblright{} (Lewis \& Short).
\end{cousinweb}

\begin{grammarlens}[title={What you've built}]
One more word for telling, explaining and describing at B1.
\end{grammarlens}

\subsection*{Your turn}
\begin{itemize}
\raggedright
  \item \emph{Say it:} \emph{siffler}
  \item \emph{Say it:} \emph{siffler}, once more
  \item \emph{Say it:} a sentence of your own with \emph{siffler}
  \item \emph{Recall:} say the French for referee, then the French for whistle, then say \emph{siffler} again
\end{itemize}

\begin{tcolorbox}[breakable,colback=teal!4,colframe=teal!35!black,title={Before you move on}]
What does siffler mean? (\textquotedblleft{}To whistle; to boo\textquotedblright{}.) Say it once more.
\end{tcolorbox}
\section[le terrain]{le terrain — pitch, ground; plot of land}

\label{lesson:FR-C505-terrain}



\begin{quote}
Before the new one: say the French for whistle, then the French for to whistle.
\end{quote}

\subsection*{You'll want to know: le terrain}
\textbf{le terrain} — \textquotedblleft{}pitch, ground; plot of land\textquotedblright{}.

\begin{quote}
Les joueurs entrent sur le terrain.
\end{quote}

\begin{quote}
\emph{The players come onto the pitch.}
\end{quote}

\emph{Un terrain de foot} — a football pitch. Also a plot: \emph{acheter un terrain}.

\begin{cousinweb}
\emph{le terrain} is inherited from Latin \emph{terrēnum} \textquotedblleft{}land, ground\textquotedblright{} (EtymDB), from \emph{terra} \textquotedblleft{}earth\textquotedblright{} (Lewis \& Short).
\end{cousinweb}

\begin{grammarlens}[title={What you've built}]
One more word for telling, explaining and describing at B1.
\end{grammarlens}

\subsection*{Your turn}
\begin{itemize}
\raggedright
  \item \emph{Say it:} \emph{le terrain}
  \item \emph{Say it:} \emph{le terrain}, once more
  \item \emph{Say it:} a sentence of your own with \emph{le terrain}
  \item \emph{Recall:} say the French for whistle, then the French for to whistle, then say \emph{le terrain} again
\end{itemize}

\begin{tcolorbox}[breakable,colback=teal!4,colframe=teal!35!black,title={Before you move on}]
What does le terrain mean? (\textquotedblleft{}Pitch, ground; plot of land\textquotedblright{}.) Say it once more.
\end{tcolorbox}
\section[marquer]{marquer — to score; to mark}

\label{lesson:FR-C505-marquer}



\begin{quote}
Before the new one: say the French for to whistle, then the French for pitch.
\end{quote}

\subsection*{You'll want to know: marquer}
\textbf{marquer} — \textquotedblleft{}to score; to mark\textquotedblright{}.

\begin{quote}
Elle a marqué deux buts dans le match.
\end{quote}

\begin{quote}
\emph{She scored two goals in the match.}
\end{quote}

\emph{Marquer un but} — to score a goal. Also to note down: \emph{marquer une adresse}.

\begin{cousinweb}
\emph{marquer} is built on \emph{marque} \textquotedblleft{}a mark, a sign\textquotedblright{}, a word of Germanic origin related to English \emph{mark}.
\end{cousinweb}

\begin{grammarlens}[title={What you've built}]
One more word for telling, explaining and describing at B1.
\end{grammarlens}

\subsection*{Your turn}
\begin{itemize}
\raggedright
  \item \emph{Say it:} \emph{marquer}
  \item \emph{Say it:} \emph{marquer}, once more
  \item \emph{Say it:} a sentence of your own with \emph{marquer}
  \item \emph{Recall:} say the French for to whistle, then the French for pitch, then say \emph{marquer} again
\end{itemize}

\begin{tcolorbox}[breakable,colback=teal!4,colframe=teal!35!black,title={Before you move on}]
What does marquer mean? (\textquotedblleft{}To score; to mark\textquotedblright{}.) Say it once more.
\end{tcolorbox}
\section[(dialogue)]{Retour sur les chapitres 498–505}

\label{lesson:FR-R505-revision-498-505}



\begin{quote}
Say the French for climbing and to climb.
\end{quote}

\subsection*{Your turn}
Say these aloud:

\begin{itemize}
\raggedright
  \item climbing, to climb, water bottle, torch, campfire
  \item board game, dice, pawn, chessboard, to cheat
  \item belote, to deal, to bet, bet, chance
  \item to collect, collector, scale model, model-making, to put together
  \item blurred, sharp, snapshot, to frame, tripod
  \item gardening, to garden, to sow, to water, watering can
  \item weight training, dumbbell, to stretch, sweat, to warm up
  \item referee, whistle, to whistle, pitch, to score
\end{itemize}

\begin{tcolorbox}[breakable,colback=teal!4,colframe=teal!35!black,title={Before you move on}]
Climbing? (\textbf{l'escalade}.) To climb? (\textbf{grimper}.) Water bottle? (\textbf{la gourde}.) Torch? (\textbf{la lampe de poche}.) Campfire? (\textbf{le feu de camp}.) Board game? (\textbf{le jeu de société}.) Dice? (\textbf{le dé}.) Pawn? (\textbf{le pion}.) Chessboard? (\textbf{l'échiquier}.) To cheat? (\textbf{tricher}.) Belote? (\textbf{la belote}.) To deal? (\textbf{distribuer}.) To bet? (\textbf{parier}.) Bet? (\textbf{le pari}.) Chance? (\textbf{le hasard}.) To collect? (\textbf{collectionner}.) Collector? (\textbf{le collectionneur}.) Scale model? (\textbf{la maquette}.) Model-making? (\textbf{le modélisme}.) To put together? (\textbf{assembler}.) Blurred? (\textbf{flou}.) Sharp? (\textbf{net}.) Snapshot? (\textbf{le cliché}.) To frame? (\textbf{cadrer}.) Tripod? (\textbf{le trépied}.) Gardening? (\textbf{le jardinage}.) To garden? (\textbf{jardiner}.) To sow? (\textbf{semer}.) To water? (\textbf{arroser}.) Watering can? (\textbf{l'arrosoir}.) Weight training? (\textbf{la musculation}.) Dumbbell? (\textbf{l'haltère}.) To stretch? (\textbf{étirer}.) Sweat? (\textbf{la sueur}.) To warm up? (\textbf{échauffer}.) Referee? (\textbf{l'arbitre}.) Whistle? (\textbf{le sifflet}.) To whistle? (\textbf{siffler}.) Pitch? (\textbf{le terrain}.) To score? (\textbf{marquer}.)
\end{tcolorbox}
